\section{总结与延伸}

本讲从一个简单但尖锐的问题出发：固定连接的神经网络如何表示每张输入都不同的 part-whole hierarchy？GLOM 的答案是，不动态分配 node hardware，而让 position-by-level grid 中的 activity 经过 recurrent interaction 形成 islands；同一 island 的近似一致向量代表一个 parse node，上下层 learned functions 表达 part-whole prediction，lateral attention 表达同层 agreement，location-conditioned neural field 让相同 whole 在不同位置生成不同 parts。

\subsection{把整讲压成六个因果步骤}

\begin{enumerate}
  \item \textbf{Dynamic state replaces dynamic memory allocation}：输入特定结构写进 activity，而不是改 weights；
  \item \textbf{Coordinate-aware prediction proposes hierarchy}：parts 向上投票 whole，whole 向下预测 parts；
  \item \textbf{Local similarity attention forms islands}：已相似的 same-level hypotheses 互相加强；
  \item \textbf{Recurrent settling delays commitment}：多步 inference 允许 ambiguity、修正与 top-down feedback；
  \item \textbf{Reconstruction and consensus train the loop}：外部 prediction error 与内部 co-distillation 共同塑造表示；
  \item \textbf{Location-conditioned decoding preserves weight sharing}：同一 whole vector 针对不同 coordinate 产生不同输出。
\end{enumerate}

\begin{importantbox}{最终教学结论}
GLOM 最值得保留的不是某个尚未验证的 update rule，而是一套 representation discipline：如果你声称模型理解对象结构，就必须说明 node 存在哪里、part-whole relation 如何编码、ambiguity 如何保留、binding 如何形成、结构如何被 loss 训练、以及哪些 intervention 能证明该结构真的被下游计算使用。
\end{importantbox}

\subsection{六个常见误区}

\begin{warningbox}{不要把设计提案写成成功史}
\begin{enumerate}
  \item 不要说 GLOM 已经解决 object-centric vision；本讲没有完整 benchmark；
  \item 不要把二维箭头当成真实 embedding dimension 或显式 pose matrix；
  \item 不要把 island 直接等同于 ground-truth segmentation mask；
  \item 不要把 consensus 等同于 truth，agreement 也可能是 shared error；
  \item 不要把 Hough analogy 当成已证明的 complexity advantage；
  \item 不要把 location-conditioned decoder 等同于完整 NeRF/modern world model。
\end{enumerate}
\end{warningbox}

\subsection{面向现代系统的谨慎连接}

今天的 object-centric representation、slot-based model、iterative refinement、equivariant network、implicit neural representation、video world model 与 sparse hierarchical attention 都触碰了 GLOM 的局部问题。真正有价值的比较不是看名字，而是逐项检查：是否有动态 entity state，是否显式处理 pose，是否保留 multiple hypotheses，是否允许 top-down correction，是否存在可读出的 part-whole relation，以及结构是否通过 causal intervention 被验证。

\teachervoice{结束前 Hinton 说这是一场复杂的 talk，最好的用途可能是鼓励大家去读长论文。这个结尾与开场的 “vaporware” 呼应：他不是宣布问题已解决，而是在邀请读者把一个 representation idea 变成更严格的模型与实验。}

\subsection{拓展阅读}

\begin{center}
\footnotesize
\begin{tabularx}{\textwidth}{>{\raggedright\arraybackslash}X>{\raggedright\arraybackslash}X}
\href{https://arxiv.org/abs/2102.12627}{Hinton, \emph{GLOM}}：完整设计文档。 &
\href{https://arxiv.org/abs/1710.09829}{Sabour et al., \emph{Capsules}}：routing 背景。 \\[0.25em]
\href{https://arxiv.org/abs/2002.05709}{Chen et al., \emph{SimCLR}}：contrastive scaffold。 &
\href{https://arxiv.org/abs/1706.03762}{Vaswani et al., \emph{Transformer}}：attention 基线。 \\[0.25em]
\href{https://arxiv.org/abs/1810.04805}{Devlin et al., \emph{BERT}}：masked reconstruction。 &
\href{https://arxiv.org/abs/1804.03235}{Anil et al., \emph{Online Distillation}}：co-distillation。 \\[0.25em]
\href{https://arxiv.org/abs/2003.08934}{Mildenhall et al., \emph{NeRF}}：coordinate-conditioned field。 & \\
\end{tabularx}
\end{center}

\end{document}
